\documentclass[../main.tex]{subfiles}
\begin{document}

\newpage

\section{Viabilidad económica}

En este apartado se estudia la viabilidad económica de sustituir el refrigerante R290 por la mezcla RE170/R1270 (80\%/20\%), ya que, esta es la que ha demostrado tener un VHC más elevado y un COP similar al del RE170. Además, su composición de propileno es más elevada, lo que abaratará su coste. Como valores de referencia, se tomarán aquellos calculados en el documento PRESUPUESTO.

Para calcular la viabilidad económica se simulará una producción de máquinas que emplean R290 como fluido de trabajo y se analizará el impacto económico que tendría la sustitución por la mezcla mencionada previamente sobre los beneficios obtenidos.

\subsection{Inversión inicial}

La inversión inicial contempla únicamente el coste del estudio teórico y experimental de la investigación, cuyo valor asciende a \textbf{9.842,15~€}.

\subsection{Explotación}

Para cuantificar tanto los beneficios como los gastos, se simulan dos escenarios: uno con un incremento del margen de beneficio y otro con un aumento de las unidades vendidas. Este planteamiento se justifica por la mejora en la eficiencia energética, que permite tanto elevar el precio del producto como incrementar su volumen de ventas. En la \autoref{tab:escenarios_viabilidad} se recogen las variables evaluadas en cada uno de los dos escenarios.

\begin{table}[ht]
\centering
\begin{tabular}{c|cc|cc|}
\cline{2-5}
 & \multicolumn{2}{c|}{\textbf{Escenario 1}} & \multicolumn{2}{c|}{\textbf{Escenario 2}} \\ \cline{2-5} 
 & \multicolumn{1}{c|}{\textbf{\begin{tabular}[c]{@{}c@{}}Margen de\\ beneficio\end{tabular}}} & \textbf{\begin{tabular}[c]{@{}c@{}}Unidades\\ anuales\\ vendidas\end{tabular}} & \multicolumn{1}{c|}{\textbf{\begin{tabular}[c]{@{}c@{}}Margen de\\ beneficio\end{tabular}}} & \textbf{\begin{tabular}[c]{@{}c@{}}Unidades\\ anuales\\ vendidas\end{tabular}} \\ \hline
\multicolumn{1}{|c|}{\textbf{Máquina con R290}} & \multicolumn{1}{c|}{30\%} & 500 & \multicolumn{1}{c|}{30\%} & 500 \\ \hline
\multicolumn{1}{|c|}{\textbf{\begin{tabular}[c]{@{}c@{}}Máquina con\\ RE170/R1270 (80\%/20\%)\end{tabular}}} & \multicolumn{1}{c|}{31\%} & 500 & \multicolumn{1}{c|}{30\%} & 550 \\ \hline
\end{tabular}
\caption{Escenarios de la viabilidad económica}
\label{tab:escenarios_viabilidad}
\end{table}

\begin{table}[ht]
\centering
\begin{tabular}{c|c|c|c|}
\cline{2-4}
 & \textbf{Coste (€)} & \textbf{Margen de beneficio} & \textbf{Precio de venta (€)} \\ \hline
\multicolumn{1}{|c|}{\textbf{Máquina con R290}} & 1.265,32 & 30\% & 1,655,91 \\ \hline
\multicolumn{1}{|c|}{\textbf{\begin{tabular}[c]{@{}c@{}}Máquina con\\ RE170/R1270 (80\%/20\%)\end{tabular}}} & 1.271,91 & 31\% & 1.666,20 \\ \hline
\end{tabular}
\caption{Precios de producción y venta de las dos máquinas}
\end{table}

Para estimar el beneficio derivado de la sustitución, se asume un período de producción de 5 años.

El beneficio bruto anual ($Bb$) se obtiene mediante la \ref{eq:bb} donde $V$ es el precio de venta y $P$ es el coste de producción de cada máquina.

\begin{equation}
\begin{aligned}
Bb_{anual} &= (V_{RE170/R1270\ (80\%/20\%)} - P_{RE170/R1270\ (80\%/20\%)}) \cdot n_{RE170/R1270 (80\%/20\%)\ anual} \\
&\quad - (V_{R290} - P_{R290}) \cdot n_{R290\ anual}
\end{aligned}
\label{eq:bb}
\end{equation}

donde $V$ es el precio de venta y $P$ es el coste de producción de cada máquina.

A partir del beneficio bruto se obtiene el beneficio neto ($Bn$) aplicando la tasa del impuesto de sociedades, fijada en el 25\% conforme a la legislación fiscal española vigente, según la expresión \ref{eq:impuestos}:

\begin{equation}
    Bn = Bb \cdot (1 - \text{impuestos})
    \label{eq:impuestos}
\end{equation}

El flujo de caja ($FC$) se calcula mediante la expresión \ref{eq:fc}. Dado que el análisis se centra exclusivamente en el incremento de beneficio generado por el cambio de refrigerante, la amortización toma un valor nulo, por lo que el flujo de caja coincide con el beneficio neto:

\begin{equation}
    FC = Bn + \text{amortización} = Bn
    \label{eq:fc}
\end{equation}

Los resultados anuales del beneficio bruto, el beneficio neto y el flujo de caja se recogen en la \autoref{tab:res_economicos}.

\begin{table}[ht]
\centering
\begin{tabular}{|c|c|c|}
\hline
\textbf{Variable (€)} & \textbf{Escenario 1} & \textbf{Escenario 2} \\ \hline
Beneficio bruto anual & 7.348,46 & 20.067,56 \\ \hline
Beneficio neto anual & 5.511,34 & 15.050,67 \\ \hline
Flujo de caja anual & 5.511,34 & 15.050,67 \\ \hline
\end{tabular}
\caption{Beneficios anuales}
\label{tab:res_economicos}
\end{table}

\subsection{VAN, TIR y PR}

Para evaluar la rentabilidad de la inversión se calculan el Valor Actual Neto (VAN), la Tasa Interna de Retorno (TIR) y el Período de Retorno (PR).

El VAN representa el valor actualizado de todos los flujos de caja generados durante el período de explotación, descontada la inversión inicial. Un VAN positivo indica que el proyecto es rentable. Se calcula mediante la \ref{eq:van}, donde el tipo de interés real ($i_{r}$) se obtiene a partir de la \ref{eq:ir}, sumando el interés bancario medio de financiación a empresas en España en 2025, fijado en el 4\%, y una prima de riesgo del 2\%, lo que resulta en un tipo de interés real del 6\%:

\begin{equation}
    VAN = -I_{0} + \sum_{n=1}^{5} \frac{FC_{n}}{(1 + i_{r})^{n}}
    \label{eq:van}
\end{equation}

\begin{equation}
    i_{r} = i_{bancario} + prima_{riesgo}
    \label{eq:ir}
\end{equation}

El TIR es la tasa de descuento que hace que el VAN sea igual a cero, es decir, el rendimiento interno del proyecto. Cuanto mayor sea el TIR respecto al tipo de interés real, más atractiva resulta la inversión. Se obtiene resolviendo la ecuación \ref{eq:tir}:

\begin{equation}
    -I_{0} + \sum_{n=1}^{5} \frac{FC_{n}}{(1 + TIR)^{n}} = 0
    \label{eq:tir}
\end{equation}

El Período de Retorno (PR) indica el tiempo necesario para recuperar la inversión inicial a partir de los flujos de caja generados. Dado que el flujo de caja es constante a lo largo de todos los años del período de explotación, se calcula mediante la expresión simplificada \ref{eq:pr}:

\begin{equation}
    PR = \frac{I_{0}}{FC_{anual}}
    \label{eq:pr}
\end{equation}

Los resultados obtenidos se recogen en la \autoref{tab:res_economicos_2}.

\begin{table}[ht]
\centering
\begin{tabular}{|c|c|c|}
\hline
\textbf{Variable} & \textbf{Escenario 1} & \textbf{Escenario 2} \\ \hline
VAN (€) & 13.373,63 & 53.556,74 \\ \hline
TIR & 48\% & 151\% \\ \hline
PR (años) & 1.79 & 0.65 \\ \hline
\end{tabular}
\caption{Resultados de la viabilidad económica}
\label{tab:res_economicos_2}
\end{table}

Los resultados obtenidos indican que el proyecto es económicamente viable: el VAN positivo confirma que la inversión genera valor, el TIR del 48\% y 151\% superan ampliamente el tipo de interés real del 6\%, lo que refleja una rentabilidad muy elevada, y el período de retorno en ambas situaciones es menor a los dos años.

\newpage